\section{性幻想}
% 移入：“性幻想”（另见第二篇“性偏好与性幻想”节，注意分工）
\subsection{什么是性幻想}

性幻想（sexual fantasy）指在清醒状态下出现的、带有性内容的想象，也可以称作“白日梦”或“意淫”——即在性交之前或进行之中，在脑海里勾勒出性行为的过程与场景。它既可以只是转瞬即逝的片段，也可以是被反复编织、相对完整的情节。白日梦不论是否带有性内容，都能使人暂时忘却现实，进入一个“什么事情都有可能发生”的空间；而带性内容的幻想之所以格外引人注目，在于它同时承担着唤起与探索的功能。

需要澄清的第一个误解是“有性幻想就是不知足或不忠”。研究显示，性幻想与关系满意度、承诺水平并无负相关——拥有稳定关系的人同样会有大量幻想，且其内容常与伴侣无关。幻想是心智的正常活动，与实际行动的意愿并不等同。

\subsection{性幻想的普遍性}

性幻想在人群中的普遍程度远超多数人的预期：多项调查显示，绝大多数成年人报告有性幻想经历，且相当比例的人在其性生活中会主动运用幻想。

\begin{itemize}
  \item \keyword{普遍性极高}：绝大多数成年人报告有性幻想，频率与内容因人差异极大。
  \item \keyword{与关系质量无关}：有幻想不代表关系有问题，也不代表不忠倾向。
  \item \keyword{内容多样}：从贴近现实的场景到完全不可能的情境，均可出现。
  \item \keyword{主动性使用}：许多人在性活动中主动借助幻想来提升唤起。
\end{itemize}

第二个需要澄清的是“幻想的内容等于真实的欲望”。研究反复显示，许多幻想包含本人不愿在现实中实现、甚至明确排斥的元素。幻想的功能往往不在“预示行为”，而在心理层面的象征性满足——例如关于“被强势引领”的幻想，未必表示想要被支配，而可能指向对“可以放下控制、不必负责”的渴望。

从发展角度看，性幻想通常自青春期开始活跃，并随经验与认知变化而调整内容。成年期的幻想内容相对稳定，但会随关系状态、生命阶段与心理状态发生幅度不等的波动。这一稳定性提示，幻想在很大程度上反映了个体相对固定的心理需求结构。

\begin{tcolorbox}[colback=pink!5!white,colframe=red!50!black,title=贴心小叮咛]
有性幻想，不代表你对现在的关系不满意，也不代表你在心里“背叛”了谁。它更像是一个只属于你自己的房间——你可以在里面放任何东西，而不用向任何人交代。
\end{tcolorbox}

\subsection{性幻想的作用}

让五花八门的做爱样式在脑海里驰骋，有助于人们在性生活中变得更有想像力、更自信，因而能获得更大的满足。性幻想还能激发并维持性欲，所以对那些性欲不容易被激发起来的人尤其有用。关于性幻想在唤起辅助、情绪调节、自我认识等方面的系统功能，以及其潜在风险（如高度固化、以幻想替代现实），另见“性偏好与性幻想”一节。

\subsection{性幻想在伴侣关系中的角色}

性幻想在伴侣关系中的作用是双向的：它可以成为增进亲密的资源，也可以成为疏离的来源。区分两者的关键，在于幻想是被分享、被接纳，还是被隐瞒、被评判。

作为资源的一面，分享幻想能带来多重收益。它使双方更了解彼此的需求结构，为探索提供方向；它本身即是一种亲密的自我暴露，能加深信任；它还可以为长期关系中的性互动注入新鲜感，对抗习惯化。研究显示，能够就幻想进行开放交流的伴侣，其性满意度通常更高。

\begin{itemize}
  \item \keyword{增进了解}：分享幻想提供关于需求的具体信息，远胜于猜测。
  \item \keyword{深化亲密}：愿意暴露私密想象，本身是一种信任的表达。
  \item \keyword{对抗习惯化}：幻想为长期关系的性互动提供新的方向与素材。
  \item \keyword{共同探索}：部分幻想可转化为双方都愿意尝试的实践。
\end{itemize}

作为风险的一面，则来自几种常见情形。第一是被迫暴露：一方在未经同意的情况下发现对方的幻想内容（如浏览记录），容易引发强烈的不安与误读。第二是评判反应：分享后遭到嘲笑、嫌恶或道德谴责，往往造成持久的信任损伤，使对方从此封闭。第三是幻想替代现实：若性活动长期依赖特定幻想才能完成，可能削弱对真实伴侣的回应能力。

因此，实践上宜遵循几条原则：分享是邀请而非要求，对方有权不参与；分享从低风险内容开始，逐步推进；明确区分“说出来”与“要去做”；以及——最重要的一条——不评判。若一方对某类幻想感到不适，可以明确表示不愿参与，但不宜以道德姿态贬低对方。安全的气氛一旦被破坏，重新打开会极为困难。

\begin{tcolorbox}[colback=pink!5!white,colframe=red!50!black,title=贴心小叮咛]
如果伴侣愿意把幻想告诉你，请把它当作一份礼物而不是一份供词。你的第一反应，很可能决定了他/她以后还会不会对你敞开心门。
\end{tcolorbox}

\subsection{常见主题与个体差异}

性幻想的主题分布广泛，但研究发现某些主题在人群中出现的频率明显更高。了解这些常见主题，有助于把个人的幻想置于正常范围之内看待，减少不必要的自我否定。

跨文化研究中报告频率较高的主题包括：与熟悉的人发生关系的想象（如朋友的伴侣、同事）、在具有情感张力的情境中（如重逢、被渴望）、权力动态相关的场景（支配或服从）、多人参与的场景、以及与自己现实伴侣的想象。这些主题在不同性别、不同文化中的分布有差异，但重叠范围很大。更细分的列举还包括：窥视他人或被他人窥视的想象、同性之间的性想象等，它们在人群中均有相当比例的报告。

\begin{itemize}
  \item \keyword{熟悉对象}：对认识但并非伴侣的人产生想象，是最常见的主题之一。
  \item \keyword{权力动态}：支配与服从的幻想普遍存在，且在不同性别中均有报告。
  \item \keyword{多人场景}：报告频率不低，但多数人并不希望现实中实现。
  \item \keyword{情境张力}：如重逢、追求、被渴望，强调情感与心理的张力而非纯粹的身体。
  \item \keyword{与伴侣的想象}：相当比例的人以现实伴侣为幻想对象，尤其在长期关系中。
  \item \keyword{窥视与暴露}：窥视他人或被他人窥视的想象，亦属常见类型。
  \item \keyword{被迫情节}：在幻想层面并不少见，多与“无需为欲望负责”的心理有关；须强调幻想与现实的非自愿行为之间没有连续性，后者是法律与伦理的禁区。
\end{itemize}

性别差异方面，统计趋势显示：男性报告的幻想频率与视觉具体性平均更高，内容更常涉及多个对象或明确的性行为场景；女性的幻想平均更常包含情感叙事、情境铺垫与关系元素。在“何时使用幻想”上也有平均层面的差异：有的人只在“前戏”中幻想，有的人不到性交开始就幻想不起来，还有的人刚结识异性时幻想纷呈、真正做爱时反倒全无；传统观察认为男性更常在性交开始前借助幻想激发性欲、开始后随即减少（此时注意力转向对射精的控制），而女性的幻想则更常贯穿全过程。需要说明的是，以上均为分布层面的平均趋势，个体重叠极大，不应据此预判具体的人。

个体差异则更为显著。研究显示，幻想内容的丰富度与个人经验、性态度、想象能力与情绪状态都相关。此外，幻想会随生命阶段变化：青春期常以探索与身份确认为主，成年期常与关系需求相交织，而老年期的幻想可能更多围绕亲密与回忆。理解这一演变，有助于把幻想视为动态的心理活动，而非固定的人格标记。

\begin{tcolorbox}[colback=pink!5!white,colframe=red!50!black,title=贴心小叮咛]
如果你曾因为自己的幻想内容而自责，不妨先知道一件事：你想到的那些，绝大多数人都想过。差别只在于谁说出来、谁没说。
\end{tcolorbox}

\subsection{性幻想与角色扮演}

很多夫妻已经习惯了在做爱中扮演平等的角色，在这种情况下尝试着玩一玩“统治”与“服从”的权力交换游戏，可能会激发起性欲来。角色扮演的关键，不在于情节本身，而在于双方是否都在知情同意的前提下、在安全的边界内参与。

\noindent\textit{【科学警示】}进行权力交换游戏时，必须严格遵循以下原则：
\begin{itemize}
  \item \textbf{SSC原则}：安全（Safe）、理智（Sane）、知情同意（Consensual）——这是所有性游戏的基石
  \item \textbf{安全词（Safe Word）}：双方约定一个明确的词语，任何一方说出这个词时，活动必须立即停止
  \item \textbf{绝对自愿}：任何一方都不应被迫做自己不愿意做的事
  \item \textbf{事后关怀（Aftercare）}：游戏结束后，双方应给予彼此情感支持和安抚
\end{itemize}

近年来，实践社群在 SSC 之外也常采用“风险知情共识”（RACK）等框架，强调在充分知情的基础上主动评估与管理风险；无论采用哪种框架，知情、自愿与可随时退出都是不可让步的底线。

\subsection{共享性幻想}

很多人喜欢将性幻想当作自己的个人秘密，而有些人恰恰相反，他们喜欢将自己的性幻想说给配偶听，还有些人甚至喜欢两个人一起编织性幻想，以便一起享受。要想共享性幻想，就需要两个人在做爱前或做爱中，像演戏一样把事先编织好的情节表演出来。只要双方愿意，这种“表演”就会成为活跃性关系、提高性兴趣的一个好办法。

要想把性幻想“表演”出来，最离不开的就是强有力的想像，同时也最好不要离开一些简单的“道具”，如软绳、遮眼带等，利用它们可以增加性生活的兴奋感。使用道具时应遵守前述的安全原则，并注意器具的清洁、材质安全与使用禁忌。

\noindent\textit{【编者注】}性幻想是个人隐私的一部分，是否与伴侣分享应完全出于自愿。重要的是：
\begin{itemize}
  \item 不要因为自己有性幻想而感到羞愧——这是正常的人类心理现象
  \item 不要强迫伴侣分享他们不想分享的幻想
  \item 不要把幻想等同于现实——幻想本身并不代表对现实关系的不满
  \item 如果选择分享，应以尊重和开放的态度进行沟通
\end{itemize}

\subsection{如何看待伴侣的性幻想}

得知伴侣的性幻想内容，对许多人而言是一次情感上的考验。尤其当幻想涉及他人、多人的场景或自己无法提供的元素时，容易引发被否定感、嫉妒与不安。如何处理这一情境，对关系影响深远。

值得先区分几个情形。第一是自己主动询问后得知，还是无意间发现。无意间发现（如看到浏览记录）往往伴随着“被隐瞒”的额外伤害，冲击通常更大。第二是幻想中是否包含具体可识别的对象。若幻想指向熟识的人，现实的嫉妒更容易被激活。第三是伴侣是否愿意分享与讨论。愿意开放沟通者，其幻想通常更接近“无害的心理活动”。

\begin{itemize}
  \item \keyword{区分性质}：先厘清是常规心理活动，还是已经涉及实际的关系外投入。
  \item \keyword{识别自身反应}：嫉妒与被否定感是正常反应，但不必立即按它行动。
  \item \keyword{避免审判}：以道德姿态谴责，只会关闭沟通，不解决任何问题。
  \item \keyword{表达感受}：可以说明自己的不适，但以“我感到”而非“你有问题”的方式。
  \item \keyword{设定边界}：就“哪些可以、哪些不行”达成明确约定，而非笼统表态。
\end{itemize}

处理的核心，是把“幻想”与“行为”分开。大量的幻想并不指向行动意图，它们的功能是心理层面的象征性满足。因此，面对伴侣的幻想，更有用的提问不是“你是不是想那么做”，而是“这个幻想对你来说意味着什么”。后者的答案通常更能揭示对方真实的心理需求，往往也更容易被接纳——例如，多人的幻想可能意味对“被强烈渴望”的向往，而非真的想要多人关系。

同时也需要承认：即便理解了机制，不适感仍然可能持续存在。这种情况下，明确表达感受并协商边界是正当的、必要的。例如，可以接受对方持有某种幻想，但明确表达不愿意被拿来做比较，或不愿在性活动中被要求扮演某个角色。把边界说清楚，比要求对方改变幻想内容更现实——后者通常做不到，也容易造成隐瞒。

\begin{tcolorbox}[colback=pink!5!white,colframe=red!50!black,title=贴心小叮咛]
幻想几乎是改不掉的，但边界是可以谈的。与其要求对方“不许想”，不如说清楚“我希望我们在现实里怎么做”——前者只会让他/她学会不说，后者才能真正解决问题。
\end{tcolorbox}
